% !TEX program = lualatex
% Prueba 3 de LuaLaTeX en TeXsting: fuentes OpenType, microtipografía y paquetes que solo existen en LuaLaTeX.
\documentclass[11pt,a4paper]{article}
\usepackage[spanish]{babel}
\usepackage{fontspec}
\setmainfont{TeX Gyre Pagella}
\setsansfont{TeX Gyre Heros}
\setmonofont{TeX Gyre Cursor}
\newfontfamily\termes{TeX Gyre Termes}
\usepackage[expansion=true,protrusion=true]{microtype}
\usepackage{lua-ul}
\usepackage{xcolor}
\usepackage{luacode}
\usepackage{lipsum}
\usepackage[margin=2.4cm]{geometry}

\title{Prueba 3: tipografía}
\author{TeXsting}
\date{\today}

\begin{document}
\maketitle

\section{Fuentes OpenType}
Texto en TeX Gyre Pagella. \textsf{Esto en Heros (sin serifa).} \texttt{Esto en Cursor (monoespaciada).}
{\termes Y esto en Termes, parecida a Times.}

\medskip
Versalitas: \textsc{Universidad Nacional de Ingeniería}.
Números de estilo antiguo: {\addfontfeatures{Numbers=OldStyle} 0123456789}
frente a los normales 0123456789.
Ligaduras: ffi, fl, ff \quad sin ligaduras: {\addfontfeatures{Ligatures=NoCommon} ffi, fl, ff}.

\section{Subrayados que cortan línea (lua-ul)}
\underLine{Este subrayado lo hace Lua y puede cruzar varias líneas sin romperse, algo que
con pdfLaTeX necesita trucos.} \strikeThrough{Texto tachado.} \highLight[yellow]{Texto resaltado
con marcador amarillo que también puede seguir en la línea siguiente sin problema.}

\section{Texto que Lua transforma}
\begin{luacode*}
  function al_reves(s)
    local t = {}
    for _, c in utf8.codes(s) do table.insert(t, 1, utf8.char(c)) end
    return table.concat(t)
  end
  function letras(p)
    tex.sprint(string.format("La palabra «%s» tiene %d letras.", p, utf8.len(p)))
  end
\end{luacode*}
«Anita lava la tina» al revés es «\directlua{tex.sprint(al_reves("Anita lava la tina"))}», y
«ñandú» es «\directlua{tex.sprint(al_reves("ñandú"))}».
\directlua{letras("electrónica")}

\section{Microtipografía}
Con \texttt{microtype} (expansión y protrusión de caracteres) los párrafos justificados
quedan más parejos:

\lipsum[1-2]

\section{Colores}
\textcolor{red!70!black}{Rojo oscuro}, \textcolor{blue!60}{azul},
\colorbox{green!15}{fondo verde claro}.
\end{document}
